\NumberedItem{2}
\HistoricTitle{\textsc{Gauss} an \textsc{Olbers}. Göttingen, 17. Oktober 1811.}
\HistoricCitation{[\textsc{Wilhelm Olbers}, Sein Leben und seine Werke II, 1, (1900), S. 482.]}
\noindent\makebox[\linewidth][s]{{.}\hfill{.}\hfill{.}\hfill{.}\hfill{.}\hfill{.}\hfill{.}\hfill{.}\hfill{.}\hfill{.}\hfill{.}\hfill{.}\hfill{.}\hfill{.}\hfill{.}\hfill{.}\hfill{.}\hfill{.}\hfill{.}\hfill{.}\hfill{.}\hfill{.}\hfill{.}\hfill{.}\hfill{.}\hfill{.}\hfill{.}}\par
\noindent\hspace*{\parindent}Meine Pallasrechnungen haben nun seit 6 Wochen ganz ruhen müssen. Ich habe mich viel diese Zeit her mit den transcendenten Functionen, worauf die Integration der Gleichung
\[
(\alpha+\beta x+\gamma xx)\frac{ddy}{dx^2}+(\delta+\varepsilon x)\frac{dy}{dx}+\zeta y=0
\]
\noindent führt, beschäftigt und sehr artige Sachen gefunden. Die meisten transcendenten Functionen, mit denen man sich bisher beschäftigt hat, sind darunter als specielle Fälle begriffen. — Auch mit den unexplicabeln Functionen
\[
1{.}\,2{.}\,3\,{.}\,{.}\,{.}\,x\ \text{und}\ 1+\frac{1}{2}+\frac{1}{3}+{.}\,{.}\,{.}+\frac{1}{x}
\]
\noindent hängt diese Untersuchung zusammen.
\par
